\documentclass[11pt]{article}

% --- packages (all standard; compiles with pdflatex) ----------------------
\usepackage[utf8]{inputenc}
\usepackage[T1]{fontenc}
\usepackage{lmodern}
\usepackage[margin=1in]{geometry}
\usepackage{booktabs}
\usepackage{amssymb}
\usepackage{microtype}
\newcommand{\cmark}{\checkmark}
\newcommand{\xmark}{\,--\,}
\usepackage{xcolor}
\usepackage{listings}
\usepackage[round]{natbib}
\usepackage[hidelinks]{hyperref}
\usepackage{enumitem}

\lstset{basicstyle=\ttfamily\small,breaklines=true,columns=fullflexible,
  frame=single,framerule=0.3pt,rulecolor=\color{gray!40},
  backgroundcolor=\color{gray!4},xleftmargin=4pt,xrightmargin=4pt}

\newcommand{\tool}{\textsc{jsoncanon}}
\newcommand{\code}[1]{\texttt{#1}}

\title{\tool{}: Trustworthy JSON Canonicalization via\\
       Dual Independent Implementations and Differential Testing}

\author{Luciano Federico Pereira\thanks{Independent researcher.
  ORCID: \href{https://orcid.org/0009-0002-4591-6568}{0009-0002-4591-6568}.
  Source, specification, and reproduction harnesses are released with this paper at
  \url{https://github.com/lucianofedericopereira/json-canon}.}}

\date{\today}

\begin{document}
\maketitle

\begin{abstract}
A hash, byte comparison, or \code{git diff} over JSON is only meaningful if the
JSON has first been reduced to a single deterministic byte stream. Many
canonicalizers exist, but their correctness usually rests on a single
implementation and a test suite written by the same author against the same
mental model. We present \tool{}, a JSON canonicalizer and adjacent toolkit
built around a different correctness discipline: \emph{two fully independent
implementations} (in Python and Nim) that are required to produce
\emph{byte-identical} output, arbitrated by a written specification and enforced
by a differential-testing gate. We argue that this dual-implementation
byte-parity contract, combined with validation against \emph{external} oracles
(a JavaScript engine for number formatting, CPython's regular-expression engine
for pattern matching), yields markedly higher confidence than single-implementation
testing, because the two implementations share no code and were written to agree
only through the specification. To our knowledge no existing tool offers this
combination: a single self-contained, dependency-free package that pairs two
byte-identical implementations with coverage of the full JSON-integrity
lifecycle---lenient ingestion of real-world dialects, canonicalization,
validation against four schema languages, deterministic binary encodings, and
structural transforms---where every nontrivial component (shortest-float
printing, SHA-256, regular-expression matching, the schema engines) is
reimplemented in-tree rather than delegated to a third-party library. The
individual capabilities exist as scattered single-language libraries; the
contribution is their integration under one specification and one correctness
contract. We describe the canonical form---in particular a
number-normalization scheme that operates purely on decimal strings and never
routes a value through a binary float---and the deliberate tension between that
lossless default and RFC~8785 (the JSON Canonicalization Scheme), which
\emph{requires} IEEE-754 rounding; we resolve the tension with an opt-in mode
that loudly reports every value whose meaning changes. We detail
determinism-by-construction techniques (generated lookup tables shared across
both languages, a shared regular-expression subset, strict per-code-point
processing) and a reproducible validation methodology (millions of randomized
inputs diffed against external oracles, dozens of cross-language parity
assertions). The result is a single deterministic byte stream over a wide
superset of real-world JSON dialects, plus interoperable validation
(JSON~Schema, JSON Type Definition, CDDL, GeoJSON, I-JSON), cross-format
conversion (CBOR, MessagePack, JData), and structural transforms (JSON Pointer,
Patch, Merge Patch). We are explicit about what the method does \emph{not}
guarantee: byte parity proves the two implementations agree, not that they are
both correct---hence the external oracles---and several schema features are
deliberate subsets.
\end{abstract}

\noindent\textbf{Keywords:} canonical JSON; deterministic serialization;
differential testing; content addressing; reproducibility; RFC~8785; shortest
float; data science; machine-learning data pipelines.

\section{Introduction}

The promise of canonical JSON is simple to state: take JSON produced by
\emph{anything}---a Python emitter, a pandas frame, a hand-edited file with
comments and trailing commas, a JavaScript service---and emit one deterministic
byte stream, so that \code{cmp}, \code{sha256}, and \code{git diff} over the
output mean something regardless of the producer. The difficulty is entirely in
the word ``deterministic.'' Producers disagree on object-member order, on
whether \code{4.0} and \code{4} are the same number, on \code{ensure\_ascii}
escaping, on whitespace, and on a long tail of dialect features. A
canonicalizer must collapse all of that to a single normal form, and---if its
output is to be used as a content address---it must do so identically across
platforms, languages, and time.

This raises a question that is rarely asked directly: \emph{how do we know a
canonicalizer is correct?} The usual answer is a test suite. But a test suite
written by the implementer encodes the implementer's understanding of the
specification; a misreading of the specification appears in both the code and
the tests, and passes. For a primitive whose entire value is that independent
parties compute the same bytes, this is an uncomfortable foundation.

\tool{} adopts a stronger discipline. It ships \emph{two} implementations, in
Python and in Nim, that share no source code. A written contract
(\code{SPEC.md}) is the sole arbiter: where the two disagree, the specification
wins and at least one implementation has a bug. A continuous gate asserts that
the two produce byte-identical output across a corpus of fixtures and a matrix
of options, and---crucially---validates the genuinely hard, non-obvious pieces
(shortest-float formatting, regular-expression matching) against \emph{external}
reference implementations rather than against each other. The thesis of this
paper is that this combination---independent reimplementation for agreement,
external oracles for correctness---is a practical and effective way to build a
trustworthy canonicalizer, and that the engineering techniques required to make
two languages agree byte-for-byte are themselves worth reporting.

The implementations total roughly 5{,}600 lines of Nim and 3{,}700 lines of
Python; the gate spans fifty-eight cross-language parity assertions and is clean under
\code{mypy~-{}-strict}, \code{pyright}, and \code{nim~check}.

\paragraph{Contributions.}
\begin{itemize}[leftmargin=1.4em,itemsep=2pt]
  \item A \emph{dual-implementation byte-parity contract} for canonical
        serialization, with a specification as arbiter and a differential gate
        as enforcement (Section~\ref{sec:arch}).
  \item A lossless number-canonicalization scheme that operates on decimal
        strings and never touches a binary float, making cross-language
        agreement independent of any float printer; and an analysis of its
        tension with RFC~8785, resolved by a loud, non-silenceable warning
        discipline (Sections~\ref{sec:spec} and~\ref{sec:jcs}).
  \item \emph{Determinism-by-construction} techniques for cross-language
        agreement: lookup tables generated once and embedded in both languages,
        a shared regular-expression subset, and strict per-code-point processing
        (Section~\ref{sec:determinism}).
  \item A reproducible validation methodology using external oracles and
        large-scale differential fuzzing (Section~\ref{sec:validation}).
  \item A \emph{self-contained, dependency-free} artifact that integrates the
        full JSON-integrity lifecycle---lenient ingestion, canonicalization, four
        schema languages, deterministic binary formats, and structural
        transforms---under a single specification and parity contract.
        We are not aware of another tool that combines these properties in one
        package (Section~\ref{sec:positioning}); the novelty is in the
        integration, the correctness discipline, and the self-containment,
        not in any individual algorithm.
\end{itemize}

We are equally explicit about the limitations (Section~\ref{sec:limitations}):
parity is agreement, not correctness; some validators are subsets; there are no
machine-checked proofs.

\section{Scope and Applications}\label{sec:scope}

\paragraph{What the tool is for (and not for).}
\tool{} is useful precisely when JSON output must be \emph{compared, hashed,
signed, deduplicated, or diffed}---that is, when the bytes, not just the parsed
value, carry meaning. It is \emph{not} a faster JSON parser, a data-processing
framework, or a database. Its design centre of gravity is integrity: the default
number engine is lossless (Section~\ref{sec:spec}), so the canonical bytes are an
exact, reversible re-presentation of the input value, suitable as a content
address. When byte-level interoperability with an external ecosystem is the
goal---most importantly JSON Web Signature and the broader JOSE
family~\cite{rfc7515}, which standardize on JCS---the opt-in JCS mode produces
exactly the bytes those systems expect, at the documented cost of binary64
rounding. The two regimes (lossless content addressing vs.\ JCS interoperability)
cover most real needs, and the tool makes the choice explicit, not implicit.

\paragraph{Data science and \texttt{pandas}.}
Tabular workflows routinely produce ``the same'' table in byte-different JSON:
columns in a different order, an integer column that became float after a
\code{groupby}, or simply a different emitter. \tool{} makes such tables
hash-equal. The package ships a \code{pandas}~\cite{pandas} accessor that
serializes a DataFrame/Series through a fixed internal orientation and then
canonicalizes, so two frames holding the same data share a SHA-256 regardless of
column order or int-vs-float dtype (Listing~\ref{lst:pandas}). This underpins
content-addressed caching of intermediate artifacts, deduplication of result
sets, and reproducible-pipeline checks where a stage is re-run only if the
canonical hash of its inputs changed.

\begin{lstlisting}[caption={Content-addressed equality of tabular data, despite
  column-order and dtype differences.},label={lst:pandas},language=Python]
import pandas as pd, jsoncanon.pandas_accessor
df1 = pd.DataFrame({"price": [4, 10],  "qty": [2, 3]})      # int64
df2 = pd.DataFrame({"qty": [2., 3.],   "price": [4., 10.]}) # float64, reordered
df1.jsoncanon.sha256() == df2.jsoncanon.sha256()            # True
\end{lstlisting}

\paragraph{Statistics and scientific computing.}
Numeric data exchange is where the lossless default earns its keep.
High-precision decimals and integers beyond $2^{53}$---survey weights,
fixed-point monetary values, large identifiers, exact rationals serialized as
decimals---are preserved character-for-character rather than silently rounded
through a float, and the two implementations agree without depending on any
platform's float printer. For array-shaped numeric data, the JData/NeuroJSON
convention~\cite{jdata} (used in neuroimaging and scientific computing) is decoded
into plain nested arrays so that N-dimensional datasets can be canonicalized,
hashed, and deduplicated like any other JSON. The result is a stable content
identity for datasets that reproduces across machines and language runtimes---a
precondition for citable, re-runnable computational results.

\paragraph{LLM and AI systems.}
Large-language-model systems generate and consume JSON constantly---tool-call
arguments, structured outputs, evaluation records---and benefit from a canonical
form in several distinct ways. \emph{(i)~Caching:} keying a prompt, tool call, or
retrieval result by the SHA-256 of its canonical form means semantically
identical payloads collide in the cache even when a model reorders object keys or
emits \code{4.0} where a previous call emitted \code{4}. \emph{(ii)~Dataset
hygiene:} training and evaluation corpora can be deduplicated, and train/test
contamination detected, by canonical-hash equality rather than fragile string
matching. \emph{(iii)~Structured-output validation:} model-produced JSON can be
checked against a JSON~Schema or JSON Type Definition before it is trusted by
downstream code, with identical verdicts in whichever language the service is
written in. \emph{(iv)~Reproducible evaluation:} logging canonical hashes of
inputs and outputs makes eval runs comparable across time and harness versions.
\emph{(v)~Integrity of identifiers:} models frequently emit large integer IDs;
the JCS-mode warning discipline (Section~\ref{sec:jcs}) surfaces exactly the cases
where naive float-based canonicalization would silently corrupt such an ID---a
failure a single-implementation tool might never reveal.

\paragraph{Security, integrity, and reproducible builds.}
Canonicalization is a prerequisite for signing and content addressing. Detached
signatures over JSON (JOSE/JWS~\cite{rfc7515}) require a canonical byte stream;
content-addressable stores, tamper-evident logs, and build/supply-chain
attestations all need ``same value $\Rightarrow$ same bytes $\Rightarrow$ same
hash.'' Here the lossless default is the safe choice: it guarantees the hashed
bytes reproduce the input exactly, and the dual-implementation discipline reduces
the risk that a verifier written in a different language computes a different
digest.

\paragraph{Configuration, DevOps, and cross-language services.}
Canonicalizing configuration before committing or comparing turns \code{git diff}
into a semantic diff (the \code{-{}-diff} mode reports changes by JSON path,
ignoring formatting and key-order noise) and makes drift between environments
detectable. For polyglot service meshes, the contract that a Python service and a
Nim service (or any future port adhering to the same specification) emit identical
canonical bytes is exactly the property needed for cross-language caches,
idempotency keys, and request signing; the CBOR and MessagePack bridges
additionally let those bytes be exchanged compactly while remaining deterministic.

\section{Background and Related Work}\label{sec:background}

\paragraph{JSON and canonicalization.}
JSON is specified by RFC~8259~\cite{rfc8259}, which fixes a grammar but
deliberately leaves member ordering and number representation to producers.
Several canonical forms have been proposed to remove that latitude. The most
prominent standardized one is the JSON Canonicalization Scheme (JCS),
RFC~8785~\cite{rfc8785}, which sorts object keys by UTF-16 code unit and
serializes numbers exactly as ECMAScript's \code{Number.prototype.toString}
would. Other ad-hoc ``canonical JSON'' conventions (e.g.\ the OLPC form) predate
JCS and differ mainly in number and escaping rules. \tool{}'s default form is
deliberately \emph{not} JCS---it preserves arbitrary-precision numbers as exact
decimals---and offers JCS as an explicit mode.

\paragraph{Shortest-float printing.}
JCS number output is the shortest decimal string that round-trips to the same
IEEE-754 binary64, with ECMAScript's exponent thresholds. The modern algorithm
for the shortest round-tripping decimal is Ryu~\cite{ryu}; we port its
double-to-string core to both languages.

\paragraph{Dialects and binary encodings.}
Real input includes JSON5~\cite{json5} (unquoted keys, comments, hex, etc.),
NDJSON, and RFC~7464 JSON text sequences~\cite{rfc7464}. Binary serializations
with their own determinism stories include CBOR~\cite{rfc8949} (whose
\S4.2 defines a core deterministic encoding) and MessagePack~\cite{msgpack}.
JData/NeuroJSON~\cite{jdata} encodes N-dimensional numeric arrays in JSON.

\paragraph{Validation.}
Schema languages relevant to JSON include JSON~Schema~\cite{jsonschema}
(Draft~2020-12 and the widely deployed Draft-07), JSON Type Definition
(RFC~8927)~\cite{rfc8927}, and CDDL (RFC~8610)~\cite{rfc8610}. RFC~7493
(I-JSON)~\cite{rfc7493} defines an interoperable subset, and RFC~7946
(GeoJSON)~\cite{rfc7946} a geospatial profile with its own structural rules.

\paragraph{Pointers and patches.}
JSON Pointer (RFC~6901)~\cite{rfc6901}, JSON Patch (RFC~6902)~\cite{rfc6902},
and JSON Merge Patch (RFC~7386)~\cite{rfc7386} address and transform documents.

\paragraph{Differential testing.}
The idea of cross-checking independent implementations to find bugs is
well established in compilers and parsers~\cite{mckeeman,csmith}. Our
contribution is not the technique but its application as a \emph{primary
correctness contract} for a canonicalizer, combined with external oracles for
the components where ``agreement'' alone would be insufficient.

\paragraph{Positioning relative to existing tools.}\label{sec:positioning}
Every capability \tool{} provides has prior art, but---to our knowledge---only
in \emph{separate}, typically single-language, typically dependency-bearing
libraries: JCS implementations exist per language; \texttt{ajv} and
\texttt{python-jsonschema} validate JSON~Schema; \texttt{cbor2} and
\texttt{msgpack} handle the binary formats; \texttt{python-jsonpatch} applies
RFC~6902. We are not aware of any single tool that simultaneously
(i)~ships \emph{two independent implementations held to byte-identical output}
as its correctness contract; (ii)~is \emph{self-contained and free of
third-party runtime dependencies}, reimplementing even shortest-float conversion
(Ryu), SHA-256, the regular-expression engine, and the schema validators
in-tree---in part precisely so that cross-language agreement does not hinge on
any external library; and (iii)~spans the \emph{entire JSON-integrity lifecycle},
from lenient ingestion of real-world dialects through canonicalization,
validation against four schema languages (JSON~Schema, JTD, and CDDL, plus the
GeoJSON and I-JSON profiles), deterministic binary re-encoding, and structural
transforms, in one artifact. The novelty claimed here is this integration and
the discipline behind it, not the rediscovery of any underlying algorithm.
Table~\ref{tab:comparison} summarizes the gap.

\begin{table}[t]
\centering
\caption{Positioning by capability category (illustrative, not an exhaustive
  feature audit; representative tools, current at time of writing).
  \cmark~present; \xmark~absent.}
\label{tab:comparison}
\small
\begin{tabular}{@{}lcccccc@{}}
\toprule
 & Dual impl. & No 3rd-party & Canonical & Schema & Binary & Struct. \\
Tool & (byte-parity) & runtime deps & form & languages & formats & transforms \\
\midrule
\tool{} (this work)      & \cmark & \cmark & \cmark & 4 & 2 & \cmark \\
JCS library (per lang.)  & \xmark & varies & \cmark (JCS) & 0 & \xmark & \xmark \\
\texttt{ajv}             & \xmark & \xmark & \xmark & 1 & \xmark & \xmark \\
\texttt{python-jsonschema} & \xmark & \xmark & \xmark & 1 & \xmark & \xmark \\
\texttt{cbor2} / \texttt{msgpack} & \xmark & varies & \xmark & 0 & 1 & \xmark \\
\texttt{json-patch}      & \xmark & \xmark & \xmark & 0 & \xmark & \cmark \\
\bottomrule
\end{tabular}
\end{table}

\section{The Canonical Form}\label{sec:spec}

The canonical output is UTF-8 with no byte-order mark, no insignificant
whitespace, separators \code{,} and \code{:}, and no trailing newline (unless
requested). Three rules carry the weight.

\paragraph{Numbers (the important one).}
Numbers are normalized \emph{as decimal strings}, using only string and
integer-digit operations; a number token is never parsed into a binary float.
Given a token, the implementation decomposes it into sign, integer digits,
fraction digits, and exponent, forms the significand, strips leading and
trailing zeros (tracking the decimal point), and re-emits a plain decimal. Thus
\code{4.0}, \code{4.50e1}, and \code{4} all canonicalize to \code{4};
\code{0.10} to \code{0.1}; \code{1e3} to \code{1000}; and a thirty-digit
integer survives exactly. Two consequences matter. First, the result is
lossless for big integers and high-precision decimals. Second, and central to
this paper, the result is \emph{independent of any language's float printer}:
because no float is involved, the Python and Nim implementations agree by
construction rather than by luck.

\paragraph{Object member order.}
Members are sorted ascending by Unicode code point---the sequence of scalar
values---which is deterministic and identical across implementations. (UTF-8
byte order coincides with code-point order, so a byte-wise sort suffices.) This
differs from JCS, which sorts by UTF-16 code unit; the two agree except for
astral-plane keys, and the JCS ordering is available in JCS mode
(Section~\ref{sec:jcs}).

\paragraph{Strings.}
Every character is emitted as raw UTF-8 except those that must be escaped:
\code{"}, \code{\textbackslash}, the five short escapes, and other C0 controls
as lowercase \code{\textbackslash u00xx}. Non-ASCII characters and \code{/} are
never escaped, which erases all \code{ensure\_ascii} differences between
producers.

The input grammar is a deliberate superset of RFC~8259: BOM autodetection;
JSON5 features (unquoted Unicode/escaped identifier keys, single quotes,
comments, trailing commas, hex and leading-plus numbers, extra string escapes,
line continuations); Python literals (\code{True}/\code{False}/\code{None});
and non-finite tokens. Every non-standard feature is normalized away on output
and, in lint mode, reported.

\section{Architecture and the Byte-Parity Contract}\label{sec:arch}

Both implementations share one internal shape: a lenient parser produces a
value tree (null, bool, number-as-raw-token, string, array, object,
non-finite), and a serializer walks the tree to produce canonical bytes.
Numbers are carried as their raw token text and normalized only at serialization
time, which is what makes the decimal-string scheme above possible.

The contract is: \emph{for the same input and options, the Python and Nim CLIs
must emit identical bytes.} \code{SPEC.md} is the arbiter; a disagreement is by
definition a bug in at least one implementation, never a matter of taste. The
contract is enforced by \code{tools/parity.sh}, which runs both implementations
over a fixture corpus and an option matrix (number formats, output encodings
with and without BOM, lint, JCS, the binary formats, the validators, the
transforms) and asserts byte-identical results against each other and against
committed golden files. A second script, \code{tools/check.sh}, wraps the parity
gate together with static analysis (\code{mypy~-{}-strict}, \code{pyright},
\code{nim~check}) and both unit-test suites. Continuous integration runs the
whole gate, plus the fuzzers of Section~\ref{sec:validation}, on every change.

This architecture turns specification ambiguity into a mechanical signal. When
the specification is silent or unclear, the two implementations tend to diverge,
the gate fails, and the ambiguity must be resolved in \code{SPEC.md} before
either implementation is ``right.''

\section{Determinism by Construction}\label{sec:determinism}

Byte parity across two languages is not free; several techniques are required
to prevent each language's own libraries from introducing drift.

\paragraph{Numbers without floats.}
As above, decimal-string normalization removes the single largest source of
cross-language disagreement---each language's float-to-string routine---from the
default path entirely.

\paragraph{Generated lookup tables.}
Where data is unavoidable, it is generated once and embedded \emph{identically}
in both languages rather than computed from each language's standard library.
Two cases: (i)~the Ryu shortest-float core needs two 128-bit power-of-five
tables (326 and 342 entries); these are produced by a generator script
(\code{tools/gen\_ryu\_tables.py}) from exact big-integer arithmetic and emitted
as a Nim source file, with the generation formulas anchored against published
constants. (ii)~JSON5 unquoted-key classification follows the ECMAScript
\code{IdentifierName} grammar, which depends on Unicode general categories;
because Python's \code{unicodedata} and Nim's \code{std/unicode} can differ by
Unicode version, we generate \code{ID\_Start}/\code{ID\_Continue} code-point
ranges once (pinned to one Unicode version) and embed the same table in both,
so classification cannot drift.

\paragraph{A shared regular-expression subset.}
JSON~Schema's \code{pattern} keyword requires regular-expression matching.
Delegating to each language's regex engine would break parity on edge cases, so
we implement one small backtracking regex engine---literals, \code{.}, character
classes, anchors, groups, alternation, the standard quantifiers, and the common
escapes---with identical semantics in both languages, and validate it
externally (Section~\ref{sec:validation}).

\paragraph{Per-code-point processing.}
Operations that index ``characters''---UTF-16 key ordering for JCS, string
length bounds for schema validation, identifier scanning---are defined on
Unicode code points, matching Python's native string indexing, and the Nim port
decodes to code points rather than operating on UTF-8 bytes where it would
differ.

\paragraph{Pure SHA-256.}
Because Nim's standard library ships no SHA-256, the \code{-{}-sha256} content
hash uses a dependency-free in-tree implementation, validated against Python's
\code{hashlib}; this keeps the digest a function of the canonical bytes alone,
with no third-party variation.

\section{The Lossless/JCS Tension}\label{sec:jcs}

The lossless decimal default and RFC~8785 are in direct conflict on numbers.
JCS mandates that each number be serialized as ECMAScript's
\code{Number.prototype.toString}: parse to IEEE-754 binary64 and print the
shortest round-tripping decimal. This necessarily rounds: the integer
\code{9007199254740993} becomes \code{9007199254740992}, and values beyond the
binary64 range are not representable at all. The lossless engine, by contrast,
would preserve \code{9007199254740993} exactly.

\tool{} treats this as a first-class design decision rather than a silent
behavior. The lossless engine is the default. JCS is an opt-in mode
(\code{-{}-jcs}) that switches both the number serializer (to a from-scratch
port of Ryu's double-to-string, followed by the ECMAScript formatting
thresholds) and the key ordering (to UTF-16 code units). When JCS rounds or
overflows a value, the tool emits a warning naming the JSON path and the exact
before/after, e.g.\ \code{9007199254740993 -> 9007199254740992 (IEEE-754
rounding; not reversible)}. The warning cannot be silently discarded: the
quiet flag that suppresses it on standard error \emph{requires} a log
destination, so a meaning-changing conversion always leaves a trace. This
``warn, never silently lose'' discipline is, we argue, the right default for a
tool whose entire purpose is that the output faithfully represents the input.

\section{Validation Methodology}\label{sec:validation}

Byte parity establishes that the two implementations \emph{agree}. It does not
establish that they are \emph{correct}: a shared misreading of the specification
could make both wrong in the same way. We address this by validating the
genuinely hard components against external reference implementations, and by
fuzzing at scale. All harnesses are in-tree and reproducible.

\paragraph{Shortest-float vs.\ a JavaScript engine.}
The Ryu port is the riskiest numeric component. We feed the bit patterns of
several million random doubles (plus hard edge cases: subnormals, the maximum
double, every exponent-threshold boundary) to the Nim implementation and compare
each output, byte for byte, to Node.js's \code{String(x)}---the very function
RFC~8785 references. We have run this to ten million doubles with zero
mismatches. Separately, a cross-language harness feeds the same number tokens to
both CLIs in JCS mode and confirms identical output, exercising the
decimal-to-double parse on both sides.

\paragraph{Regex subset vs.\ CPython.}
The shared regex engine is validated two ways simultaneously over tens of
thousands of randomly generated pattern/text pairs: each result is compared
against CPython's \code{re.search} (correctness for the supported subset) and
against the Nim port (cross-language parity). Both held at forty thousand pairs.

\paragraph{Half-precision floats, exhaustively.}
The CBOR encoder selects the shortest of float16/32/64. The float16 codec is
verified \emph{exhaustively} against Python's \code{struct} over all $2^{16}$
bit patterns.

\paragraph{Salvage mode, by mutation.}
The best-effort recovery mode (\code{-{}-force}) is fuzzed over thousands of
randomly corrupted documents, asserting that both implementations produce
identical recovered output \emph{and} identical recovery reports. This fuzzer
surfaced a latent parity bug---malformed number tokens that the lenient parser
accepted but the serializer later rejected with differing error messages---which
we fixed by validating number tokens at parse time in both languages.

\paragraph{The committed gate.}
Beyond fuzzing, the parity gate asserts fifty-eight byte-identical outcomes across
fixtures and the option/format/validator/transform matrix, and both unit-test
suites pass under strict static analysis. Table~\ref{tab:validation} summarizes.

\begin{table}[t]
\centering
\caption{Validation harnesses and their oracles.}
\label{tab:validation}
\small
\begin{tabular}{@{}lll@{}}
\toprule
Component & Oracle & Scale \\
\midrule
Ryu shortest float & Node.js \code{String(x)} & $10^7$ random doubles \\
JCS number tokens & cross-impl (Nim==Python) & $5\times10^5$ tokens \\
Regex subset & CPython \code{re} \emph{and} cross-impl & $4\times10^4$ pairs \\
CBOR float16 codec & Python \code{struct} & all $2^{16}$ patterns \\
\code{-{}-force} salvage & cross-impl (output + report) & $4\times10^3$ mutations \\
Full canonical form & cross-impl + golden files & 58 parity assertions \\
\bottomrule
\end{tabular}
\end{table}

\section{Feature Surface}\label{sec:features}

On top of canonicalization, the same value-tree substrate supports a range of
JSON-adjacent operations, each implemented twice and held to the same parity
contract. Table~\ref{tab:features} lists them. Three categories are worth
noting. \emph{Validators} report conformance issues at a JSON path:
I-JSON~(RFC~7493), GeoJSON~(RFC~7946), JTD~(RFC~8927), CDDL~(RFC~8610, a
practical subset), and JSON~Schema (Draft~2020-12 and Draft-07, including
\code{\$ref}/\code{\$anchor}, the applicators, \code{unevaluated*}, and opt-in
\code{format} assertion). \emph{Cross-format} conversion uses deterministic
encodings: CBOR core deterministic encoding (RFC~8949~\S4.2) and an analogous
deterministic MessagePack, both round-tripped and byte-checked across
implementations. \emph{Transforms} run before canonicalization: JSON Pointer
extraction, JSON Patch, and JSON Merge Patch.

\begin{table}[t]
\centering
\caption{Capabilities, all dual-implemented and parity-checked.}
\label{tab:features}
\small
\begin{tabular}{@{}lll@{}}
\toprule
Area & Capability & Reference \\
\midrule
Canonical form & lossless decimal numbers, code-point key sort & this work \\
                & JCS mode (Ryu numbers, UTF-16 sort) & RFC 8785 \\
Input dialects  & JSON5 keys/escapes/numbers & JSON5 \\
                & NDJSON; JSON text sequences; concatenated & RFC 7464 \\
                & \code{-{}-force} salvage of malformed input & this work \\
Cross-format    & CBOR \code{-{}-from}/\code{-{}-to} & RFC 8949 \\
                & MessagePack \code{-{}-from}/\code{-{}-to} & MessagePack \\
                & JData N-D arrays \code{-{}-from jdata} & NeuroJSON \\
Validation      & I-JSON; GeoJSON & RFC 7493/7946 \\
                & JTD; CDDL (subset) & RFC 8927/8610 \\
                & JSON Schema 2020-12/draft-07 (+\code{-{}-format}) & JSON Schema \\
Transforms      & JSON Pointer; Patch; Merge Patch & RFC 6901/6902/7386 \\
Utility         & lint; \code{-{}-check}; \code{-{}-sha256}; \code{-{}-diff} & this work \\
                & output re-encoding (UTF-8/16/32, latin-1, BOM) & this work \\
\bottomrule
\end{tabular}
\end{table}

\section{Limitations}\label{sec:limitations}

We state the boundaries of the method plainly.

\begin{itemize}[leftmargin=1.4em,itemsep=2pt]
  \item \textbf{Parity is agreement, not correctness.} Two implementations can
        share a misreading of the specification and agree on the wrong output.
        This is precisely why the hard components are validated against
        \emph{external} oracles (Node.js, CPython's \code{re}, Python
        \code{struct}); for components without an external oracle, confidence
        rests on the specification and the test suite, as in any project.
  \item \textbf{No machine-checked proofs.} Correctness is established
        empirically (differential fuzzing) and by review, not by formal
        verification.
  \item \textbf{Deliberate subsets.} The CDDL validator covers common
        constructs, not the full grammar (no generics, sockets, or control
        operators); JSON~Schema omits remote \code{\$ref} and \code{\$dynamicRef}
        (which require network resolution / dynamic scope) and the
        \code{content*} keywords, and its \code{unevaluated*} support covers the
        common applicators rather than the full recursive annotation algorithm.
        The bundled regex engine is a subset of ECMA-262.
  \item \textbf{In-memory model.} Documents are parsed into a value tree;
        multi-gigabyte single documents are out of scope (record-oriented
        NDJSON and JSON text sequences are already processed incrementally).
  \item \textbf{Lossy by request.} JCS, CBOR, and MessagePack route numbers
        through binary64 (and integers beyond 64 bits, for MessagePack, through
        a float); this is inherent to those formats and is reported rather than
        hidden.
\end{itemize}

\section{Conclusion}\label{sec:conclusion}

\tool{} is, to our knowledge, the first self-contained, dependency-free toolkit
to span the full JSON-integrity lifecycle in a single artifact, and it
demonstrates that a canonicalization tool can be built around a correctness
discipline stronger than single-implementation testing: two
independent implementations forced to agree byte-for-byte through a written
specification, with the genuinely hard components cross-checked against external
reference implementations. The engineering required to make two languages agree
exactly---decimal-only number normalization, generated lookup tables shared
across languages, a shared regex subset, strict per-code-point processing---is
reusable beyond this project. The same substrate cleanly supports a broad
surface of JSON-adjacent functionality (dialect ingestion, validation against
four schema languages, deterministic binary formats, and structural transforms)
without compromising the parity contract. We believe the dual-implementation,
externally-oracled approach is a good fit for any primitive whose value depends
on independent parties computing the same bytes.

\paragraph{Availability and reproduction.}
The two implementations, the specification, the golden fixtures, the validation
harnesses, and the continuous-integration configuration are released together at
\url{https://github.com/lucianofedericopereira/json-canon}. A
single script, \code{tools/reproduce.sh}, rebuilds both implementations and runs
the full gate followed by every differential fuzzer (each with a fixed seed)
against its external oracle, reproducing Table~\ref{tab:validation} end to end;
\code{tools/reproduce.sh quick} runs the same pipeline at reduced scale for a
fast check. Every quantitative claim in this paper is reproducible by that one
command.

\bibliographystyle{abbrvnat}
\bibliography{references}

\end{document}
